% page 108 du fac-simile (image p-107 du scan)
% bloc 162.6 x 233.3 mm ; marge gauche 8.0 mm
\pgc{-12.700mm}{0.000mm}{
\l{\cel{3}lOO}
\l{}
\l{\sou{dekano}.\cel{2}I. Alta-rangiero di kirko (en kirko katedrala, la}
\l{\cel{3}prezidero di la kanonikaro). II. Alta-rangiero di univers-}
\l{\cel{3}itato (la profesoro selektita kom lua chefo dum ula quanto}
\l{\cel{3}de tempo). - DEFIRS.}
\l{}
\l{\sou{dekantar}.\cel{2}(\sou{trans.}) Transvarsar (liquido) tale ke en la vazo}
\l{\cel{3}unesma restez la lizo. - DEFIS.}
\l{}
\l{\sou{dekapar}. (\sou{trans.}) Deprenar la strato de oxido qua kovras la}
\l{\cel{3}surfaco di metalo. - DFIS.}
\l{}
\l{\sou{dekastero}.\cel{2}lO steri.}
\l{}
\l{\sou{deklamar}. (\sou{trans.})\cel{2}Recitar, saliigante la senco per la toni}
\l{\cel{3}e la gesti. - DEFIRS.}
\l{}
\l{\sou{deklarar}. (\sou{trans.})\cel{2}I. Konocigar aperte (ulo). - II. Enuncar}
\l{\cel{3}aperte fakto. - DEFIRS.}
\l{}
\l{\sou{deklinar}. (\sou{trans.}) I. Enuncar la serio de la dezinenci (di}
\l{\cel{3}kazi, di genri, di nombri) tra qua pasas nomo, pronomo,}
\l{\cel{3}adjektivo qualifikala en ula lingui. - II. (\sou{pri astro}).}
\l{\cel{3}Distar de la equatoro di la sfero cielala, supere o sube.}
\l{\cel{3}III. (\sou{pri la magnet-agulo}). Deviacar de la meridiano ter-}
\l{\cel{3}globala, per ula angulo. - DEFIS.}
\l{}
\l{\sou{dekoktar}. (\sou{trans}.) Boliigar substanco en liquido, por ke la}
\l{\cel{3}utilajo quan lu kontenas dissolvesez en ta liquido. - DEFIRS.}
\l{}
\l{\sou{dekoltar}. (\sou{trans.})\cel{2}Abatar addope o sesgar vesto por videbl-}
\l{\cel{3}igar la kolo, e c. - DEFR.}
\l{}
\l{\sou{dekorar}. (\sou{trans.}) Garnisar ye acesoraji destinita a plubel-}
\l{\cel{3}igar. - DEFIRS.}
\l{}
\l{\sou{dekremento}. (\sou{radiotekniko}). Deskresko di la ocili o di elektro-}
\l{\cel{3}fluo.}
\l{}
\l{\sou{dekremetro}. (\sou{radiotekniko}).Aparato uzata por mezurar nemediate}
\l{\cel{3}la dekremento logaritmala di cirkuito de elektro-fluo}
\l{\cel{3}alternanta o di ondo-serio elektromagnetala.}
\l{}
\l{\sou{dekretar}. (\sou{trans.}) (\sou{aludante sutoritatozo suprega}). Decida,}
\l{\cel{3}rezolvar. - DEFIRS.}
\l{}
\l{\sou{dekubitar}. (\sou{netrans}.) (\sou{aludante homo-korpo, kushita)}. Pozesar}
\l{\cel{3}(flankale, dorsale). - DEFIS.}
\l{}
\l{\sou{delegar}. (\sou{trans., ad ulu)}. Charjar (ulu) per ofico, transmis-}
\l{\cel{3}ante a lu omna, o parto de, sua autoritato. - DEFIRS.}
\l{}
\l{\sou{delektar}. (\sou{trans}.) Igar (ulu) savurar juo. - DEFIS.}
}
